% Expanded front matter for the current Density4Drug draft.
% Style target: short claim-led paragraphs and compact sentences, following
% the four reference papers supplied for the 260910 writing pass.
% Citation keys follow the existing 260903 drafts and must be synchronized
% with the BibTeX database used by the current manuscript.

% =============================================================================
% ABSTRACT
% =============================================================================

\begin{abstract}
Protein machine learning depends on representations that can transfer across
tasks. Most current representations use amino-acid sequences or fitted atomic
coordinates. They omit the experimental evidence used to determine many of
those coordinates. Electron-density maps can retain local support and
conformational evidence that a single fitted model does not expose. We ask
whether this signal can provide a reusable representation for protein--ligand
modeling. We introduce \ours{}, a
coordinate-and-density voxel encoder pre-trained by masked reconstruction on
more than 100K curated protein--ligand instances. \ours{} jointly encodes ligand
atoms, pocket atoms, experimental $2mF_o-DF_c$ density, and its gradient
magnitude. We transfer frozen pooled or spatial features to binding-affinity
regression and reference-assisted molecular generation. On
leakage-controlled affinity benchmarks, \ours{} improves over a matched
coordinate-only encoder and obtains the strongest point estimates on a clean
92-complex CASF-2016 subset. When fused with VoxBind, it improves mean Vina Dock
by $0.51$~kcal/mol and raises the high-affinity rate by $6.4$ percentage points.
Mean ECFP4 similarity to the reference ligand does not increase, although
PoseCheck reveals a remaining trade-off in generated-pose quality. These
results support experimental density as a transferable complement to fitted
coordinates.
\end{abstract}

% =============================================================================
% INTRODUCTION
% =============================================================================

\section{Introduction}
\label{sec:intro}

Reusable representations are central to modern protein machine learning.
Protein language models learn evolutionary and functional patterns from
amino-acid sequences \citep{rives2021biological,lin2023evolutionary}.
Structure predictors instead recover accurate atomic models from sequence
\citep{jumper2021highly,abramson2024accurate}. These representations now support
function prediction, molecular docking, affinity estimation, and drug design
\citep{radivojac2013large,ozturk2018deepdta,pinheiro2024structure}. Their success
also makes the information retained by the input representation increasingly
important.

Fitted coordinates are an incomplete record of crystallographic evidence. In
X-ray crystallography, an atomic model is refined against a diffraction-derived
electron-density map \citep{adams2011phenix}. The final coordinates provide one
concise interpretation of that map, but they do not retain its full local
signal. Weak density can indicate alternative conformations, while local map
support varies across atoms and molecular fragments
\citep{lang2010automated,lang2014protein,meyder2017estimating}. Such information
is especially relevant near a ligand, where small geometric changes can alter
intermolecular contacts.

Experimental density has already shown value in several protein--ligand tasks.
Prior work uses it for ligand recognition, local model assessment, affinity
prediction, and molecular generation
\citep{kowiel2019automatic,miyaguchi2021machine,suriana2025beyond,wang2022pocket}.
These studies establish that the map contains useful task-level information.
They do not yet show whether density-pre-trained representations can support
distinct predictive and generative tasks. This leaves representation transfer,
rather than access to density itself, as the central question.

We introduce \ours{}, a transferable coordinate-and-density encoder for
protein--ligand complexes. Each complex is represented on a common voxel grid
using ligand atoms, pocket atoms, experimental density, and density-gradient
magnitude. We curate more than 100K crystallographic instances from PLINDER and
pre-train a Channel-ViT by masked reconstruction
\citep{durairaj2024plinder,bao2023channel-24a,he2022masked}. The objective asks
the model to recover missing atomic and density regions from their shared
context. At transfer time, the pretrained encoder remains frozen and only
lightweight task-specific modules are learned.

We evaluate \ours{} in two complementary transfer settings. Binding-affinity
regression tests whether pooled features capture interaction strength. We use
complex-level leakage controls, sequence-similarity filters, and an exact-PDB
held-out CASF-2016 subset. Reference-assisted molecular generation tests whether
spatial features improve ligand construction inside a target pocket. We assess
docking, molecular properties, reference-ligand similarity, and pose geometry
so that a better interaction score is not mistaken for uniformly better
molecules.

The experiments show that experimental density provides useful information
beyond a matched coordinate-only representation. \ours{} improves affinity
prediction across the main LP-PDBBind evaluations and achieves the strongest
point estimates on the clean CASF-2016 subset. Its spatial features also improve
VoxBind's Vina scores without increasing mean ECFP4 similarity to the reference
ligand. At the same time, PoseCheck identifies a remaining gap in generated-pose
quality relative to VoxBind. This trade-off locates both the value of the density
signal and an important target for future transfer methods.

Our contributions are summarized as follows:
\begin{itemize}
    \item We introduce \ours{} and a large-scale curation pipeline for joint
    representation learning from fitted coordinates and experimental
    electron-density maps.
    \item We evaluate the frozen representation for binding-affinity regression
    under complex-level leakage controls, protein-sequence-similarity filters,
    and a clean CASF-2016 held-out set.
    \item We transfer frozen spatial features to reference-assisted molecular
    generation and evaluate predicted interactions, molecular quality,
    reference similarity, and pose geometry.
\end{itemize}

% =============================================================================
% RELATED WORK
% =============================================================================

\section{Related Work}
\label{sec:related_works}

\textbf{Experimental electron-density maps.}
An electron-density map retains experimental information that is not explicit
in a fitted coordinate set. A standard maximum-likelihood-weighted map combines
observed amplitudes with fitted-model terms using coefficients
$(2mF_o-DF_c)e^{\mathrm{i}\phi_c}$ \citep{read1986improved}. Although the map
remains partly model-dependent, its observed amplitudes can reveal weak or
heterogeneous local structure
\citep{lang2010automated,lang2014protein}. Map-based methods have been developed
for ligand identification and local structure assessment
\citep{kowiel2019automatic,miyaguchi2021machine}. We use the map as a source of
pre-training supervision and as part of a reusable complex representation.

\textbf{Protein--ligand representation learning.}
Protein--ligand models differ in both their input resolution and their training
objective. Sequence methods operate on protein residues and ligand strings,
whereas structure methods encode atoms, coordinates, distances, or voxel grids.
Recent self-supervised approaches pre-train complex or pocket encoders through
geometric denoising, masked modeling, or cross-molecular interaction prediction
\citep{jin2023dsmbind,feng2024protein,gao2024self}. These methods show that
unlabeled structures can provide transferable interaction features. \ours{}
extends this direction by aligning coordinate channels with experimental density
during pre-training.

\textbf{Voxel and masked representation learning.}
Voxel grids provide a shared spatial frame for molecular signals with different
physical meanings. They have supported both affinity prediction and generative
modeling, but their cost grows quickly with spatial resolution
\citep{stepniewska2018development,pinheiro2024structure}. Masked autoencoding
reduces this redundancy by learning to reconstruct missing regions from visible
context \citep{he2022masked}. Channel-aware Transformers further preserve the
identity of distinct input groups while allowing information exchange between
them \citep{bao2023channel-24a}. These properties make masked, channel-aware
voxel modeling a natural basis for combining atoms with experimental maps.

\textbf{Binding-affinity regression.}
Binding-affinity regression tests whether a representation captures the
strength of a protein--ligand interaction. Sequence-based models use protein
sequences and ligand SMILES without explicit three-dimensional geometry
\citep{ozturk2018deepdta,huang2021moltrans,shenoy2026nesso1}. Structure-based
models instead use voxel CNNs, molecular graphs, or equivariant networks
\citep{stepniewska2018development,kong2026hbgsa,kong2023generalist}. Evaluation
is sensitive to overlap between training and test complexes, which motivates
leak-aware splits and similarity-stratified analyses
\citep{li2026leak,wei2026honestaffinity}. Most closely related,
\citet{suriana2025beyond} compares atom-based and experimental-density voxel
inputs; we focus on large-scale pre-training and frozen transfer under these
controls.

\textbf{Structure-based drug design.}
Pocket-conditioned generators construct molecules directly within a target
binding site. Autoregressive methods place atoms or fragments sequentially
\citep{luo20213d,peng2022pocket2mol}, while diffusion models jointly denoise
atom types and coordinates \citep{guan20233d-5e4,schneuing2024structure}.
Voxel-based models instead generate continuous atom-density grids
\citep{pinheiro20233d-81b,pinheiro2024structure}. Reference information can be
introduced through decomposed ligand priors or experimental density
\citep{guan2024decompdiff,wang2022pocket,li2025electron}. We differ by
transferring frozen spatial features learned from coordinates and experimental
density into an existing voxel generator.
